\documentclass[article]{jss}

%% -- LaTeX packages and custom commands ---------------------------------------
\usepackage{amsmath,amssymb,amsfonts}
\usepackage{algorithmic}
\usepackage{booktabs}

%% -- Article metainformation --------------------------------------------------

\author{David Puelz\\University of Austin \And
        Panos Toulis\\University of Chicago \And
        Shuxiao Huang\\University of Chicago}
\Plainauthor{David Puelz, Panos Toulis, Shuxiao Huang}

\title{\pkg{BicliqueRT}: Randomization Tests for Causal Effects Under General Interference in \proglang{R}}
\Plaintitle{BicliqueRT: Randomization Tests for Causal Effects Under General Interference in R}
\Shorttitle{\pkg{BicliqueRT}: Randomization Tests Under Interference}

\Abstract{
We present \pkg{BicliqueRT}, an \proglang{R} package for conducting exact,
finite-sample valid randomization tests of causal effects in experiments
with interference.  In such settings, the treatment assigned to one unit
may affect the outcomes of others, invalidating standard randomization
tests that assume no interference.  The package implements the
biclique-based conditional randomization test framework of
\citet{puelz2022graph}, which constructs a null-exposure graph from the
experimental design and null hypothesis, decomposes it into bicliques, and
conditions on a biclique containing the observed treatment allocation to
obtain a valid test.  \pkg{BicliqueRT} supports arbitrary interference
structures---including spatial, network, and clustered
interference---through a flexible interface for specifying exposure
mappings and null hypotheses.  Two decomposition algorithms are available:
the Bimax algorithm for binary inclusion-maximal biclusters and a greedy
heuristic.  We illustrate the package with examples covering spatial
networks, clustered experiments, and contrast hypotheses with confidence
intervals.
}

\Keywords{causal inference, interference, randomization test, biclique decomposition, \proglang{R}}
\Plainkeywords{causal inference, interference, randomization test, biclique decomposition, R}

\Address{
  David Puelz\\
  University of Austin\\
  Austin, TX, United States of America\\
  E-mail: \email{dpuelz@uaustin.org}\\

  Panos Toulis\\
  Booth School of Business\\
  University of Chicago\\
  Chicago, IL, United States of America\\

  Shuxiao Huang\\
  Department of Statistics\\
  University of Chicago\\
  Chicago, IL, United States of America
}

\begin{document}


%% -- Introduction -------------------------------------------------------------

\section{Introduction} \label{sec:intro}

Randomized experiments are the gold standard for estimating causal effects.
In the classical framework pioneered by \citet{fisher1935design} and
formalized by \citet{rubin1980randomization}, the potential outcomes of
each unit are assumed to depend only on that unit's own treatment
assignment.  This assumption, known as the \emph{stable unit treatment
value assumption} (SUTVA) or \emph{no interference}, is violated in many
modern experimental settings.  For example, in social network experiments,
vaccinating one individual may protect their neighbors; in spatial
experiments, treating an agricultural plot may affect adjacent plots; and
in clustered experiments, treating one household member may spill over to
other members.

The presence of interference complicates causal inference because the
potential outcomes of each unit can depend on the entire vector of treatment
assignments.  A growing literature has developed methods for causal
inference under interference, including work by
\citet{hudgens2008toward}, \citet{sobel2006randomized},
\citet{aronow2017estimating}, and \citet{basse2019randomization}.

\citet{puelz2022graph} proposed a graph-theory approach to
randomization tests under general interference.  Their key insight is to
construct a \emph{null-exposure graph} that encodes which units and
treatment assignments are compatible under the null hypothesis.  A
biclique decomposition of this graph yields conditional randomization
tests that are valid in finite samples under arbitrary interference
patterns.

The \pkg{BicliqueRT} package implements this methodology in \proglang{R}
\citep{R}.  It provides a complete workflow for:
\begin{enumerate}
  \item Specifying an experiment design, exposure mapping, and null
        hypothesis;
  \item Constructing and decomposing the null-exposure graph using the
        Bimax algorithm \citep{prelic2006systematic} or a greedy
        heuristic;
  \item Conducting conditional randomization tests with user-specified or
        built-in test statistics;
  \item Computing confidence intervals by inverting the test.
\end{enumerate}

The remainder of this paper is organized as follows.
Section~\ref{sec:method} reviews the methodological background.
Section~\ref{sec:package} describes the package structure and core
functions. Section~\ref{sec:illustrations} provides worked examples.
Section~\ref{sec:comparison} compares \pkg{BicliqueRT} with related
software. Section~\ref{sec:summary} concludes.


%% -- Methodological Background ------------------------------------------------

\section{Methodological background} \label{sec:method}

\subsection{Setup and notation}

Consider an experiment with $N$ units indexed by $i = 1, \ldots, N$.  Let
$Z = (Z_1, \ldots, Z_N) \in \{0,1\}^N$ denote the vector of treatment
assignments, where $Z_i = 1$ indicates that unit $i$ is treated.  Under
interference, the potential outcome of unit $i$ under the full treatment
vector $z$ is $Y_i(z)$.  The observed outcome is $Y_i^{\text{obs}} =
Y_i(Z)$.

\subsection{Exposure mappings and null hypotheses}

An \emph{exposure mapping} is a function $f_i: \{0,1\}^N \to
\mathcal{E}$ that maps the full treatment vector to an exposure value for
unit~$i$.  For example, $f_i(z) = z_i$ corresponds to no interference
(the exposure of $i$ is simply its own treatment), while $f_i(z) = (z_i,
\sum_{j \in \mathcal{N}_i} z_j)$ captures both direct treatment and the
number of treated neighbors.

A null hypothesis is specified through the exposure mapping and an
equivalence relation.  For instance, the null of no interference states
\[
  H_0: Y_i(z) = Y_i(z') \quad \text{for all } i \text{ and all } z, z'
  \text{ such that } f_i(z) = f_i(z').
\]

\subsection{The null-exposure graph}

Given the exposure mapping, null hypothesis, and a set of treatment
assignments $\mathcal{Z} = \{z^{(1)}, \ldots, z^{(M)}\}$ drawn from the
experimental design, the \emph{null-exposure graph} is a bipartite graph
$G = (\mathcal{U}, \mathcal{Z}, E)$ where $\mathcal{U} = \{1, \ldots,
N\}$ is the set of units and $\mathcal{Z}$ is the set of assignments.  An
edge $(i, z^{(m)})$ exists if and only if unit $i$'s exposure under
$z^{(m)}$ is equivalent to its exposure under the observed assignment
$z^{\text{obs}}$, as determined by the null hypothesis.

\subsection{Biclique decomposition and conditional tests}

A \emph{biclique} in $G$ is a complete bipartite subgraph: a set of units
$\mathcal{F} \subseteq \mathcal{U}$ (the \emph{focal units}) and a set
of assignments $\mathcal{A} \subseteq \mathcal{Z}$ (the \emph{focal
assignments}) such that every unit--assignment pair is connected.  Under
the null hypothesis, the potential outcomes of all focal units are
invariant across all focal assignments.

The key result of \citet{puelz2022graph} is that conditioning on a
biclique yields a valid randomization test:
\begin{enumerate}
  \item Decompose the null-exposure graph into bicliques that partition
        the assignment space.
  \item Identify the biclique containing the observed assignment
        $z^{\text{obs}}$.
  \item Compute a test statistic $T(Y, z)$ for each assignment $z$ in the
        biclique.
  \item The p-value is the proportion of assignments in the biclique that
        yield a test statistic at least as extreme as the observed value.
\end{enumerate}

This conditional test controls the Type~I error rate at level $\alpha$ in
finite samples, regardless of the form of interference, provided the null
hypothesis is correctly specified.

\subsection{Decomposition algorithms}

\pkg{BicliqueRT} implements two decomposition algorithms:

\paragraph{Bimax.}  The Bimax algorithm \citep{prelic2006systematic} finds
binary inclusion-maximal biclusters.  It uses a divide-and-conquer
strategy: pick a column, divide rows by whether they have a 1 or 0 in
that column, and recurse.  The algorithm is controlled by minimum row
(\code{minr}) and column (\code{minc}) size parameters.

\paragraph{Greedy.}  A greedy heuristic that iteratively selects random
units and adds them to the focal set as long as the resulting biclique
has at least \code{mina} assignments.  This approach is faster but may
yield less balanced bicliques.


%% -- The Package --------------------------------------------------------------

\section[The BicliqueRT package]{The \pkg{BicliqueRT} package} \label{sec:package}

\subsection{Installation}

The package can be installed from CRAN:
\begin{Code}
install.packages("BicliqueRT")
\end{Code}
or from GitHub:
\begin{Code}
remotes::install_github("dpuelz/BicliqueRT")
\end{Code}

\subsection{Core workflow}

The \pkg{BicliqueRT} package supports two main workflows, corresponding
to two classes of null hypotheses.

\subsubsection{Generalized interface}

For arbitrary null hypotheses, the workflow is:
\begin{enumerate}
  \item Define the hypothesis via a list with three components:
        \code{design\_fn} (the randomization mechanism), \code{exposure\_i}
        (the exposure mapping), and \code{null\_equiv} (the equivalence
        relation under the null).
  \item Call \code{biclique.decompose()} to construct the null-exposure
        graph and decompose it into bicliques.
  \item Call \code{clique\_test()} with a test statistic to conduct the
        conditional randomization test.
\end{enumerate}

\begin{Code}
hypothesis <- list(
  design_fn  = function() rbinom(N, 1, 0.3),
  exposure_i = function(z, i) z[i],
  null_equiv = function(e1, e2) identical(e1, e2)
)
controls <- list(method = "greedy", mina = 10, num_randomizations = 2000)
decomp <- biclique.decompose(Z, hypothesis, controls)
tstat  <- gen_tstat(G, type = "htn")
result <- clique_test(Y, Z, tstat, decomp)
\end{Code}

\subsubsection{Contrast interface}

For contrast-type null hypotheses $H_0: Y_i(b) = Y_i(a) + \tau$ for all
$i$, the dedicated function \code{clique\_test\_contrast()} handles the
entire workflow:

\begin{Code}
result <- clique_test_contrast(
  Y = Y, Z = Z, Z_a = Z_a, Z_b = Z_b,
  Zobs_id = 1, alpha = 0.05, tau = 0,
  decom = "greedy", mina = 10
)
\end{Code}

\subsection{Key functions}

Table~\ref{tab:functions} summarizes the main functions.

\begin{table}[t!]
\centering
\begin{tabular}{ll}
\toprule
\textbf{Function} & \textbf{Description} \\
\midrule
\code{biclique.decompose()} & Decompose the null-exposure graph \\
\code{clique\_test()}       & Generalized randomization test \\
\code{clique\_test\_contrast()} & Contrast hypothesis test \\
\code{gen\_tstat()}         & Generate built-in test statistics \\
\code{one\_sided\_test()}   & One-sided p-value \\
\code{two\_sided\_test()}   & Two-sided p-value \\
\code{clique\_diagnostic()} & Diagnostic for tuning parameters \\
\code{CI\_grid()}           & Confidence interval (grid method) \\
\code{CI\_bisection()}      & Confidence interval (bisection method) \\
\bottomrule
\end{tabular}
\caption{\label{tab:functions}Main functions in \pkg{BicliqueRT}.}
\end{table}

\subsection{Test statistics}

The \code{gen\_tstat()} function provides three test statistics from
\citet{athey2018exact}:
\begin{itemize}
  \item \code{"elc"}: Edge-level contrast, measuring the difference in
        average outcomes between focal units with treated vs.\ untreated
        neighbors.
  \item \code{"score"}: A score-type statistic based on the correlation
        between residuals and peer treatment exposure.
  \item \code{"htn"}: Has-treated-neighbor statistic, measuring the
        correlation between outcomes and having at least one treated
        neighbor.
\end{itemize}

Users can also supply custom test statistics as functions with signature
\code{function(y, z, is\_focal)}.

\subsection{Confidence intervals}

Confidence intervals are obtained by inverting the test.  For a contrast
hypothesis $H_0(\tau): Y_i(b) = Y_i(a) + \tau$, the $1-\alpha$
confidence interval is the set of $\tau$ values for which $H_0(\tau)$ is
not rejected at level $\alpha$.  Two methods are available:
\begin{itemize}
  \item \code{CI\_grid()}: Evaluates the test on a grid of $\tau$ values.
  \item \code{CI\_bisection()}: Uses bisection to find the interval
        endpoints more efficiently.
\end{itemize}

\subsection{Parallel computation}

The \pkg{foreach} backend \citep{pkg:foreach} is used for parallelizing
the grid-based confidence interval computation.  Users can register a
parallel backend (e.g., via \pkg{doParallel}) before calling the
test.


%% -- Illustrations ------------------------------------------------------------

\section{Illustrations} \label{sec:illustrations}

\subsection{Spatial interference}

We demonstrate the generalized interface by testing the null hypothesis of
no interference in a synthetic spatial network.

\begin{Code}
library("BicliqueRT")
set.seed(42)
num_units <- 50

# Generate a spatial network (3 Gaussian clusters)
net <- out_example_network(num_units)
D <- net$D

# Adjacency matrix: connect units within distance 0.15
G <- (D < 0.15 & D > 0) * 1

# Treatment and outcomes
Z <- rbinom(num_units, 1, 0.3)
Y <- rnorm(num_units) + 0.5 * Z

# Hypothesis: no interference (exposure = own treatment)
hypothesis <- list(
  design_fn  = function() rbinom(num_units, 1, 0.3),
  exposure_i = function(z, i) z[i],
  null_equiv = function(e1, e2) identical(e1, e2)
)

controls <- list(method = "greedy", mina = 10, num_randomizations = 500)
decomp <- biclique.decompose(Z, hypothesis, controls)

tstat <- gen_tstat(G, type = "htn")
result <- clique_test(Y, Z, tstat, decomp)
result$p.value
\end{Code}

\subsection{Clustered interference}

For clustered experiments, we use the contrast interface with the
\citet{basse2018model} data-generating process.

\begin{Code}
set.seed(123)
N <- 90; K <- 30; numrand <- 200

Zprime <- out_Zprime(N, K, equal = TRUE, numrand = numrand)
Z_a <- (Zprime == 1) * 1   # spillover
Z_b <- (Zprime == 0) * 1   # pure control
Z_treat <- (Zprime == 2) * 1

Zobs_id <- 1
bf <- out_bassefeller(N, K, Zobs = Zprime[, Zobs_id], tau_main = 0)

result <- clique_test_contrast(
  Y = bf$Yobs, Z = Z_treat, Z_a = Z_a, Z_b = Z_b,
  Zobs_id = Zobs_id, decom = "greedy", mina = 10
)
result$decision
\end{Code}

\subsection{Network interference with extent-of-interference null}

Here we test a more complex null hypothesis: the extent-of-interference
null, where the exposure includes both own treatment and the count of
treated neighbors.

\begin{Code}
set.seed(99)
num_units <- 40

# Erdos-Renyi graph
G <- matrix(0, num_units, num_units)
for (i in 1:(num_units - 1))
  for (j in (i + 1):num_units)
    if (runif(1) < 0.1) G[i, j] <- G[j, i] <- 1

Z <- rbinom(num_units, 1, 0.3)
Y <- rnorm(num_units) + 0.3 * (G %*% Z)

hypothesis <- list(
  design_fn  = function() rbinom(num_units, 1, 0.3),
  exposure_i = function(z, i) c(z[i], sum(G[i, ] * z)),
  null_equiv = function(e1, e2) identical(e1, e2)
)

controls <- list(method = "greedy", mina = 5, num_randomizations = 300)
decomp <- biclique.decompose(Z, hypothesis, controls)

tstat <- gen_tstat(G, type = "score")
result <- clique_test(Y, Z, tstat, decomp)
result$p.value
\end{Code}

\subsection{Confidence intervals for contrast hypotheses}

Confidence intervals are obtained by inverting the contrast test over a
grid of $\tau$ values.

\begin{Code}
set.seed(123)
N <- 90; K <- 30; numrand <- 200

Zprime <- out_Zprime(N, K, equal = TRUE, numrand = numrand)
Z_a <- (Zprime == 1) * 1
Z_b <- (Zprime == 0) * 1
Z_treat <- (Zprime == 2) * 1

Zobs_id <- 1
bf <- out_bassefeller(N, K, Zobs = Zprime[, Zobs_id], tau_main = 0.5)

result_ci <- clique_test_contrast(
  Y = bf$Yobs, Z = Z_treat, Z_a = Z_a, Z_b = Z_b,
  Zobs_id = Zobs_id, decom = "greedy", mina = 10,
  ret_ci = TRUE, ci_method = "grid", ci_dec = 1
)
result_ci$ci
\end{Code}


%% -- Comparison with related software -----------------------------------------

\section{Comparison with related software} \label{sec:comparison}

Several \proglang{R} packages address randomization inference, but none
handle general interference.  Table~\ref{tab:comparison} summarizes the
comparison.

\begin{table}[t!]
\centering
\begin{tabular}{lccc}
\toprule
\textbf{Package} & \textbf{Interference} & \textbf{Finite-sample valid} & \textbf{General networks} \\
\midrule
\pkg{ri2} \citep{pkg:ri2}          & No & Yes & No \\
\pkg{coin} \citep{pkg:coin}        & No & Yes & No \\
\pkg{interferenceCI} \citep{pkg:interferenceCI} & Partial & Yes & No \\
\pkg{BicliqueRT}                    & General & Yes & Yes \\
\bottomrule
\end{tabular}
\caption{\label{tab:comparison}Comparison of \proglang{R} packages for
randomization inference.}
\end{table}

\paragraph{\pkg{ri2}.} The \pkg{ri2} package \citep{pkg:ri2} implements
randomization inference for standard randomized experiments, assuming no
interference.  It provides sharp null tests, confidence intervals, and
handles various randomization designs.  However, it does not account for
interference.

\paragraph{\pkg{coin}.} The \pkg{coin} package \citep{pkg:coin}
implements a general framework for conditional inference via permutation
tests.  While flexible, it assumes exchangeability of units under the
null, which is violated under interference.

\paragraph{\pkg{interferenceCI}.} The \pkg{interferenceCI} package
\citep{pkg:interferenceCI} provides confidence intervals under partial
interference, where units are grouped into clusters and interference
occurs only within clusters.  This is a special case of the general
interference handled by \pkg{BicliqueRT}.

\paragraph{Advantages of \pkg{BicliqueRT}.}  The key advantage is
generality: \pkg{BicliqueRT} handles arbitrary interference structures
specified through user-defined exposure mappings, without assuming partial
interference or any particular network structure.  The tests are exact
(finite-sample valid) and the framework naturally extends to contrast
hypotheses and confidence intervals.

\paragraph{Limitations.}  The computational cost of the biclique
decomposition grows with the number of units and randomizations.  For
very large experiments ($N > 1000$), the greedy algorithm is recommended
over Bimax.  Additionally, the power of the test depends on the quality
of the biclique decomposition; the \code{clique\_diagnostic()} function
helps users choose tuning parameters.


%% -- Summary ------------------------------------------------------------------

\section{Summary and future work} \label{sec:summary}

The \pkg{BicliqueRT} package provides \proglang{R} users with tools for
conducting exact randomization tests under general interference.  The
package implements the methodology of \citet{puelz2022graph}, supporting
both generalized and contrast-type null hypotheses, two decomposition
algorithms, and built-in test statistics for network experiments.

Future directions include:
\begin{itemize}
  \item Improved decomposition algorithms with better worst-case
        guarantees.
  \item Extension to multi-valued treatments and continuous exposures.
  \item Integration with design-based estimators for average treatment
        effects under interference.
  \item Faster C-level implementation of the null-exposure graph
        construction.
\end{itemize}

The package is available on CRAN at
\url{https://CRAN.R-project.org/package=BicliqueRT} and on GitHub at
\url{https://github.com/dpuelz/BicliqueRT}.


%% -- Acknowledgments ----------------------------------------------------------

\section*{Acknowledgments}

We thank Guillaume Basse and Avi Feller for foundational contributions
to the methodology.


%% -- Bibliography -------------------------------------------------------------

\bibliography{references}

\end{document}
